\documentclass[12pt,a4paper]{report}

% ================= PAKET =================
\usepackage[indonesian]{babel}
\usepackage[utf8]{inputenc}
\usepackage[T1]{fontenc}
\usepackage{times}
\usepackage[a4paper,left=4cm,top=4cm,right=3cm,bottom=3cm]{geometry}
\usepackage{setspace}
\usepackage{graphicx}
\usepackage{svg}
\usepackage{xcolor}
\usepackage{listings}
\usepackage{longtable}
\usepackage{array}
\usepackage{enumitem}
\usepackage{hyperref}
\usepackage{caption}
\usepackage{float} % Paket untuk posisi gambar presisi [H]

\hypersetup{
    colorlinks=true,
    linkcolor=black,
    citecolor=black,
    urlcolor=blue,
    bookmarksopen=true,
    pdftitle={Laporan Proyek Coffee E-Nose - Firmware Akuisisi dan OTA Update},
}

\onehalfspacing

% Optimasi penempatan gambar agar tidak membuat halaman kosong
\renewcommand{\topfraction}{0.85}
\renewcommand{\bottomfraction}{0.8}
\renewcommand{\textfraction}{0.15}
\renewcommand{\floatpagefraction}{0.7}

% ================= KONFIGURASI LISTING KODE =================
\definecolor{codebg}{rgb}{0.96,0.96,0.96}
\definecolor{codekw}{rgb}{0.0,0.0,0.6}
\definecolor{codecom}{rgb}{0.35,0.35,0.35}
\definecolor{codestr}{rgb}{0.0,0.45,0.25}

\lstdefinelanguage{Rust}{
    keywords={fn, let, mut, if, else, match, loop, while, for, in, struct, enum, impl, pub, use, async, await, return, continue, break, const, static, true, false, Some, None, Ok, Err, self, Self, crate, mod},
    keywordstyle=\color{codekw}\bfseries,
    sensitive=true,
    comment=[l]{//},
    morecomment=[s]{/*}{*/},
    string=[b]",
}

\lstset{
    backgroundcolor=\color{codebg},
    basicstyle=\ttfamily\footnotesize,
    breaklines=true,
    breakatwhitespace=false,
    frame=single,
    rulecolor=\color{black!30},
    numbers=left,
    numberstyle=\tiny\color{black!50},
    numbersep=8pt,
    keywordstyle=\color{codekw}\bfseries,
    commentstyle=\color{codecom}\itshape,
    stringstyle=\color{codestr},
    showstringspaces=false,
    captionpos=b,
    tabsize=4,
    columns=fullflexible,
}

% ================= DOKUMEN =================
\begin{document}

% ---------- HALAMAN SAMPUL ----------
\begin{titlepage}
    \centering
    \singlespacing
    \vspace*{0.5cm}

    {\bfseries LAPORAN EVALUASI TENGAH SEMESTER (ETS)\\PROJECT-BASED LEARNING (PBL)\par}
    \vspace{4pt}
    {Mata Kuliah: Teknologi IoT\par}

    \vspace{1.8cm}

    {\bfseries\Large SISTEM AKUISISI DATA \textit{ELECTRONIC NOSE} (\textit{E-NOSE}) DAN\\
    OTA (\textit{OVER-THE-AIR}) \textit{FIRMWARE UPDATE} BERBASIS ESP32-S3\\
    DAN THINGSBOARD \textit{CLOUD}\par}

    \vspace{0.8cm}

    % Logo resmi ITS
    \includegraphics[width=3cm]{logoits.png}
    \vspace{0.8cm}

    {Dosen Pengampu:\\
    \bfseries Ahmad Radhy, S.Si., M.Si\par}

    \vspace{0.8cm}

    {Disusun Oleh Kelompok 7:\par}
    \vspace{4pt}

    \begin{tabular}{ll}
        \bfseries Faris Ahmad Holili & 2042241028 \\[6pt]
        \bfseries Yusuf Febri Andrian & 2042241043 \\[6pt]
        \bfseries Mohammad Eka Jauhar Ismail & 2042241054 \\[6pt]
        \bfseries Reksa Putra Armediansyah & 2042241126 \\
    \end{tabular}

    \vfill

    {\bfseries DEPARTEMEN TEKNIK INSTRUMENTASI\\
    FAKULTAS VOKASI\\
    INSTITUT TEKNOLOGI SEPULUH NOPEMBER\\
    SURABAYA\\
    2026\par}
    \vspace*{0.5cm}

\end{titlepage}

% ---------- DAFTAR ISI ----------
\pagenumbering{roman}
\tableofcontents
\newpage
\pagenumbering{arabic}

% =======================================================
\chapter{Pendahuluan}

\section{Latar Belakang}
Kopi merupakan komoditas dengan variasi jenis dan profil aroma yang sangat beragam,
sehingga identifikasi jenis/varietas kopi secara objektif menjadi topik yang relevan
untuk diteliti. Salah satu pendekatan yang dapat digunakan adalah \textit{electronic
nose} (\textit{E-Nose}), yaitu sistem berbasis \textit{array} sensor gas yang meniru cara kerja indera
penciuman manusia untuk mengenali profil aroma suatu zat.

Proyek \textbf{Coffee \textit{E-Nose}} ini mengembangkan sistem akuisisi data berbasis
mikrokontroler ESP32-S3 yang membaca sinyal dari delapan kanal sensor gas melalui dua
buah ADC eksternal ADS1115, serta sensor suhu dan kelembapan DHT22 sebagai data
lingkungan pendukung. Data yang diakuisisi dikirim secara \textit{real-time} melalui
protokol MQTT ke platform IoT \textit{cloud} \textbf{ThingsBoard}, sekaligus dicatat secara lokal
melalui antarmuka serial untuk keperluan pencatatan data mentah.

Karena perangkat direncanakan digunakan pada beberapa sesi pengambilan data yang
terpisah, kebutuhan untuk dapat memperbarui \textit{firmware device} dari jarak jauh tanpa harus
membongkar/mencolok ulang kabel USB menjadi penting. Oleh karena itu, sistem ini juga
mengimplementasikan mekanisme \textbf{OTA (\textit{Over-The-Air}) \textit{firmware update}} yang
memanfaatkan fitur bawaan ThingsBoard untuk mendistribusikan berkas \textit{firmware} baru ke
\textit{device} secara nirkabel.

\section{Tujuan}
\begin{enumerate}[label=\arabic*.]
    \item Mengimplementasikan \textit{firmware} akuisisi data sensor gas \textit{array} (8 kanal) dan
    sensor lingkungan (DHT22) pada mikrokontroler ESP32-S3 menggunakan bahasa Rust.
    \item Mengimplementasikan konektivitas WiFi dan protokol MQTT untuk mengirimkan data
    telemetri secara \textit{real-time} ke ThingsBoard \textit{Cloud}.
    \item Mengimplementasikan dan menguji mekanisme OTA \textit{firmware update} \textit{end-to-end},
    mulai dari proses \textit{assign firmware} di ThingsBoard hingga \textit{device} berhasil
    mengunduh, memverifikasi, dan menjalankan \textit{firmware} baru.
    \item Mendokumentasikan kendala teknis yang ditemukan selama pengujian OTA di
    perangkat keras (\textit{hardware}) nyata beserta solusi yang diterapkan.
\end{enumerate}

\section{Ruang Lingkup}
Laporan ini berfokus pada pengembangan sisi \textbf{\textit{firmware}} (\textit{embedded software}) dan
mekanisme \textbf{OTA \textit{update}}-nya. \textit{Pipeline} pengolahan data lanjutan (ekstraksi fitur,
pelatihan model klasifikasi \textit{machine learning}) berada di luar cakupan laporan ini dan
ditangani secara terpisah pada folder \texttt{04\_features/} dan \texttt{05\_models/}
dalam struktur repositori proyek.

% =======================================================
\chapter{Dasar Teori}

\section{\textit{Electronic Nose} dan \textit{Array} Sensor Gas}
Sistem \textit{E-Nose} pada proyek ini menggunakan delapan kanal sensor gas semikonduktor \textit{metal oxide}
(MOS), yaitu MQ-3, MQ-6, MQ-7, MQ-135, TGS2600, TGS2602, TGS2611, dan TGS2620.
Setiap sensor memiliki sensitivitas berbeda terhadap senyawa volatil tertentu (misalnya
alkohol, gas mudah terbakar, senyawa organik volatil), sehingga kombinasi keluaran
kedelapan sensor membentuk suatu \textit{fingerprint} aroma yang dapat dibedakan antar
sampel kopi.

Karena keluaran sensor berupa sinyal analog, dibutuhkan ADC (\textit{Analog-to-Digital
Converter}) beresolusi tinggi. Digunakan dua modul \textbf{ADS1115} (ADC 16-bit,
antarmuka I2C, 4 kanal per modul) dengan alamat I2C berbeda (\texttt{0x48} dan
\texttt{0x49}) sehingga total 8 kanal sensor dapat dibaca melalui satu \textit{bus} I2C yang
sama (SDA pada GPIO8, SCL pada GPIO9). Sensor DHT22 digunakan untuk mencatat suhu dan
kelembapan ruangan sebagai variabel kontrol, terhubung pada GPIO10.

\section{ESP32-S3 dan \textit{Firmware} Berbasis Rust}
ESP32-S3 dipilih sebagai mikrokontroler utama karena memiliki konektivitas WiFi
bawaan, memori \textit{buffer} yang cukup untuk komunikasi MQTT, serta dukungan dua slot
partisi OTA. \textit{Firmware} ditulis dalam bahasa \textbf{Rust} menggunakan \textit{crate}
\texttt{esp-hal} (\textit{Hardware Abstraction Layer} resmi untuk keluarga \textit{chip} Espressif) dan
\textit{runtime} asinkron \texttt{embassy}, yang memungkinkan beberapa \textit{task}
(pembacaan sensor, komunikasi MQTT, dan lain-lain) berjalan secara konkuren tanpa RTOS
(\textit{Real-Time Operating System}) tradisional.

\section{Protokol MQTT dan Platform ThingsBoard}
MQTT (\textit{Message Queuing Telemetry Transport}) adalah protokol \textit{publish-subscribe}
ringan yang umum digunakan pada sistem IoT. \textbf{ThingsBoard} adalah platform IoT
\textit{cloud} yang menyediakan \textit{broker} MQTT, \textit{dashboard} visualisasi, penyimpanan
\textit{attribute} (\textit{client-side} dan \textit{shared/server-side}), serta fitur manajemen OTA
\textit{firmware update} bawaan yang digunakan langsung pada proyek ini tanpa perlu membangun
\textit{server} OTA sendiri.

\section{Mekanisme OTA (\textit{Over-The-Air}) \textit{Firmware Update}}
Konsep OTA \textit{update} pada ESP32 (dan \textit{chip} Espressif pada umumnya) mengandalkan skema
\textbf{\textit{dual-partition}}: terdapat dua partisi aplikasi (\texttt{ota\_0} dan
\texttt{ota\_1}) berukuran sama, dan \textit{device} hanya menjalankan salah satunya pada satu
waktu. \textit{Firmware} baru selalu ditulis ke partisi yang sedang tidak aktif,
sehingga apabila proses unduh gagal di tengah jalan, partisi yang sedang berjalan
tidak tersentuh sama sekali dan \textit{device} tetap dapat berfungsi normal. Setelah data
\textit{firmware} baru selesai diunduh dan diverifikasi (biasanya melalui \textit{checksum}
CRC32), \textit{bootloader} diarahkan untuk mem-\textit{boot} dari partisi baru
tersebut pada \textit{reboot} berikutnya.

Pada proyek ini, mekanisme \textit{low-level} penulisan ke partisi OTA ditangani oleh
\textit{crate} Rust \texttt{esp-hal-ota}, sedangkan protokol pertukaran data \textit{firmware}
antara \textit{device} dan ThingsBoard mengikuti protokol resmi ThingsBoard berbasis topik MQTT
\texttt{v1/devices/me/attributes} (untuk metadata versi) dan
\texttt{v2/fw/request/\{id\}/chunk/\{n\}} (untuk mengunduh \textit{firmware} per-\textit{chunk}).

% =======================================================
\chapter{Metodologi dan Implementasi Sistem}

\section{Arsitektur Sistem}
Secara umum, arsitektur sistem terdiri atas tiga lapisan:
\begin{enumerate}[label=\arabic*.]
    \item \textbf{Lapisan akuisisi} -- ESP32-S3 membaca 8 kanal sensor gas via dua
    ADS1115 serta satu sensor DHT22, kemudian mencetak data ke antarmuka serial dalam
    format CSV setiap 1 detik.
    \item \textbf{Lapisan konektivitas} -- ESP32-S3 terhubung ke jaringan WiFi lokal,
    kemudian membentuk koneksi MQTT ke \texttt{thingsboard.cloud} untuk mengirim
    telemetri dan menerima perintah/\textit{attribute} (termasuk metadata OTA).
    \item \textbf{Lapisan manajemen \textit{firmware} (OTA)} -- ThingsBoard menyimpan paket
    \textit{firmware} pada \textit{Packages repository}, yang dapat di-\textit{assign} ke
    \textit{device} tertentu. \textit{Device} secara otomatis mendeteksi penetapan (\textit{assignment})
    tersebut dan mengunduhnya bila versinya berbeda dari yang sedang berjalan.
\end{enumerate}

\section{Perangkat Keras (\textit{Hardware})}
\begin{itemize}
    \item Mikrokontroler: ESP32-S3 (\textit{flash} 16 MB)
    \item ADC eksternal: 2x ADS1115 (I2C, alamat \texttt{0x48} dan \texttt{0x49})
    \item Sensor gas: MQ-3, MQ-6, MQ-7, MQ-135, TGS2600, TGS2602, TGS2611, TGS2620
    (total 8 kanal, terbagi merata pada 2 modul ADS1115)
    \item Sensor lingkungan: DHT22 (suhu \& kelembapan), pada GPIO10
    \item \textit{Bus} I2C: SDA = GPIO8, SCL = GPIO9
\end{itemize}

\section{Struktur Partisi \textit{Flash}}
Agar mendukung OTA, tabel partisi \textit{flash} dirancang dengan dua slot aplikasi seperti
diperlihatkan pada Tabel~\ref{tab:partisi}.

\begin{table}[H]
\centering
\caption{Tabel partisi \textit{flash} ESP32-S3 (\texttt{partitions.csv})}
\label{tab:partisi}
\begin{tabular}{|l|l|l|l|l|}
\hline
\textbf{Nama} & \textbf{Tipe} & \textbf{Subtipe} & \textbf{Offset} & \textbf{Ukuran} \\
\hline
nvs      & data & nvs   & 0x9000  & 24 KB \\
otadata  & data & ota   & 0xF000  & 8 KB \\
phy\_init & data & phy  & 0x11000 & 4 KB \\
ota\_0   & app  & ota\_0 & 0x20000 & 1.5 MB \\
ota\_1   & app  & ota\_1 & 0x1A0000 & 1.5 MB \\
\hline
\end{tabular}
\end{table}

Partisi \texttt{ota\_0} dan \texttt{ota\_1} masing-masing berukuran 1,5 MB, cukup
untuk menampung ukuran biner \textit{firmware} hasil kompilasi (kurang lebih 600 KB pada
proyek ini), dengan ruang cadangan untuk pengembangan fitur berikutnya.

\section{Tahapan Pengembangan}
Pengembangan \textit{firmware} dilakukan secara bertahap dan iteratif berdasarkan riwayat
\textit{commit} pada repositori proyek. Tahapan utamanya adalah sebagai berikut.

\subsection{Fase 1 -- Inisialisasi Proyek}
Struktur folder proyek, \textit{skeleton firmware}, serta konfigurasi awal WiFi dan MQTT
disiapkan. Konfigurasi kredensial (SSID, \textit{password}, \textit{host} ThingsBoard, dan token
akses) awalnya dicoba melalui \textit{environment variable}, kemudian dipindahkan ke berkas
teks \texttt{wifi\_config.txt} yang di-\textit{gitignore} agar kredensial tidak
ikut ter-\textit{commit} ke repositori publik.

\subsection{Fase 2 -- Protokol Akuisisi 8 Sensor}
Mengikuti instruksi dosen pengampu, protokol akuisisi diperbarui menjadi 8 kanal
sensor gas melalui dua modul ADS1115, dengan 10 kali percobaan pengambilan data per
sesi. \textit{Bug} kritis pada komunikasi UART yang menyebabkan \textit{firmware hang} total
juga ditemukan dan diperbaiki pada fase ini, disertai penambahan pencatatan
(\textit{logging}) lokal sebagai cadangan apabila koneksi \textit{cloud} terputus.

\subsection{Fase 3 -- \textit{Rewrite} Akuisisi Sensor Murni}
\textit{Firmware} ditulis ulang agar berfokus murni pada akuisisi sensor sesuai pemetaan \textit{pin}
dan format keluaran yang ditetapkan pada dokumen instruksi resmi, termasuk perbaikan
\textit{pin} DHT22 (dari asumsi awal GPIO4 menjadi GPIO10 sesuai referensi repositori acuan)
dan penambahan pemindaian (\textit{scan}) \textit{bus} I2C saat \textit{boot} untuk memverifikasi
keberadaan perangkat pada \textit{bus}.

\subsection{Fase 4 -- Penggabungan Ulang Konektivitas dan OTA}
Setelah modul akuisisi sensor stabil, fungsi WiFi/MQTT digabungkan kembali ke dalam
\textit{firmware} yang sama, disusul dengan integrasi ulang mekanisme OTA \textit{firmware update}
menggunakan \textit{crate} \texttt{esp-hal-ota} beserta tabel partisi khusus OTA.

\subsection{Fase 5 -- \textit{Debugging} OTA \textit{End-to-End} di Perangkat Keras}
Fase ini merupakan fase pengujian paling intensif, mencakup penemuan dan perbaikan
tiga kendala signifikan yang dibahas secara rinci pada Bab~\ref{ch:pengujian}.

\section{Implementasi \textit{Firmware} Akuisisi Sensor}
\textit{Firmware} membaca kedelapan kanal sensor gas dan data DHT22 setiap 1 detik, lalu
mencetak hasilnya ke antarmuka serial dalam format CSV agar mudah diproses lebih
lanjut:

\begin{lstlisting}[language=Rust,caption={Format keluaran data akuisisi sensor}]
DATA,t_s,mq3,mq6,mq7,mq135,tgs2600,tgs2602,tgs2611,tgs2620,temp_c,rh_pct,dht_age_s
\end{lstlisting}

Nilai setiap kanal gas dikonversi langsung ke satuan tegangan (volt) berdasarkan LSB
ADS1115 pada mode \textit{full-scale range} $\pm4{,}096$ V, sehingga keluaran lebih
mudah diinterpretasikan tanpa perlu konversi manual di sisi pengolahan data.

\section{Implementasi Konektivitas WiFi dan MQTT}
Setelah koneksi WiFi berhasil dan alamat IP diperoleh melalui DHCP, \textit{firmware}
membentuk koneksi MQTT ke \texttt{thingsboard.cloud:1883} menggunakan token akses
\textit{device} sebagai kredensial (\textit{username}). Setelah terkoneksi, \textit{firmware}:
\begin{enumerate}[label=(\alph*)]
    \item Men-\textit{subscribe} topik \texttt{v1/devices/me/attributes} untuk
    menerima pembaruan \textit{attribute} secara \textit{real-time}.
    \item Melaporkan versi \textit{firmware} yang sedang berjalan (\texttt{current\_fw\_title},
    \texttt{current\_fw\_version}) sebagai \textit{client attribute}.
    \item Meminta (\textit{request}) nilai \textit{shared attribute} terkini terkait
    OTA (\texttt{fw\_title}, \texttt{fw\_version}, \texttt{fw\_size},
    \texttt{fw\_checksum}, \texttt{fw\_checksum\_algorithm}) melalui topik
    \texttt{v1/devices/me/attributes/request/1}.
    \item Mempublikasikan data telemetri sensor setiap 5 detik ke topik
    \texttt{v1/devices/me/telemetry}, disertai mekanisme \textit{ping} berkala agar
    koneksi MQTT tidak terputus akibat \textit{keep-alive timeout}.
\end{enumerate}

Apabila koneksi MQTT terputus, \textit{firmware} secara otomatis mencoba menghubungkan
kembali (\textit{reconnect}) setelah jeda 3 detik, sementara \textit{task} akuisisi sensor dan pencatatan
lokal tetap berjalan tanpa terganggu.

\section{Implementasi OTA \textit{Firmware Update}}
Alur lengkap mekanisme OTA yang diimplementasikan adalah sebagai berikut:

\begin{enumerate}[label=\arabic*.]
    \item \textit{Device} melaporkan versi \textit{firmware} yang sedang berjalan sebagai
    \textit{client attribute} (\texttt{current\_fw\_title} / \texttt{current\_fw\_version}).
    \item \textit{Device} meminta nilai \textit{shared attribute} terkait \textit{firmware} yang
    ditetapkan (\textit{assigned}) untuknya di ThingsBoard.
    \item \textit{Device} membandingkan versi yang ditetapkan dengan versi yang sedang
    berjalan. Apabila sama, proses dihentikan (\textit{skip}); apabila berbeda,
    proses unduh dimulai.
    \item \textit{Firmware} diunduh secara bertahap per-\textit{chunk} berukuran 4096 \textit{byte}
    melalui topik \texttt{v2/fw/request/\{id\}/chunk/\{n\}}, dan setiap \textit{chunk}
    yang diterima langsung ditulis ke partisi OTA yang tidak aktif menggunakan
    \textit{crate} \texttt{esp-hal-ota}.
    \item Nilai CRC32 dihitung secara independen atas seluruh data yang diunduh, lalu
    dibandingkan dengan nilai \textit{checksum} yang dikirim ThingsBoard.
    \item Apabila \textit{checksum} cocok, data \textit{firmware} di-\textit{flush} secara
    final ke \textit{flash} dan \textit{device} melakukan \textit{restart} untuk berpindah menjalankan
    \textit{firmware} baru dari partisi yang baru saja ditulis.
\end{enumerate}

Potongan kode berikut menunjukkan bagian pengecekan versi yang ditambahkan untuk
mencegah pengunduhan ulang \textit{firmware} yang sama (dibahas lebih lanjut pada
Bab~\ref{ch:pengujian}):

\begin{lstlisting}[language=Rust,caption={Pengecekan versi sebelum memulai unduhan OTA (\texttt{main.rs})}]
if let Some(update) = parse_fw_update(text) {
    if update.title == CURRENT_FW_TITLE
        && update.version == CURRENT_FW_VERSION
    {
        info!(
            "OTA: {} v{} sudah versi yang jalan sekarang, skip download",
            update.title, update.version
        );
        continue;
    }

    info!(
        "OTA package terdeteksi: {} v{} ({} bytes)",
        update.title, update.version, update.size
    );
    // ... proses unduh chunk demi chunk ...
}
\end{lstlisting}

Implementasi lengkap struktur program \textit{firmware} pada berkas \texttt{main.rs} diperlihatkan pada Gambar~\ref{fig:ss_main_rs}.

\begin{figure}[H]
    \centering
    \includegraphics[width=0.75\textwidth]{main rs.png}
    \caption{Tangkapan layar kode \textit{firmware} utama (\texttt{main.rs}) berbasis Rust dan Embassy.}
    \label{fig:ss_main_rs}
\end{figure}

% =======================================================
\chapter{Pengujian dan Hasil}
\label{ch:pengujian}

Pengujian OTA dilakukan langsung pada perangkat keras (\textit{hardware}, bukan simulasi), dengan cara
meng-\textit{assign} paket \textit{firmware} baru ke \textit{device} melalui \textit{dashboard} ThingsBoard dan
mengamati proses unduhan melalui \textit{serial monitor} (\texttt{espflash monitor}).
Pada iterasi pengujian pertama, ditemukan tiga kendala signifikan yang harus
diselesaikan sebelum OTA dapat berfungsi secara andal. Ringkasannya diperlihatkan
pada Tabel~\ref{tab:kendala}.

\begin{table}[H]
\centering
\caption{Ringkasan kendala OTA yang ditemukan dan solusinya}
\label{tab:kendala}
\begin{tabular}{|p{0.5cm}|p{3.2cm}|p{4.3cm}|p{4.3cm}|}
\hline
\textbf{No} & \textbf{Gejala} & \textbf{Penyebab} & \textbf{Solusi} \\
\hline
1 &
Verifikasi gagal / \textit{checksum} CRC32 tidak cocok &
ThingsBoard mengirim nilai \textit{checksum} \texttt{fw\_checksum} dengan urutan
\textit{byte} (\textit{endianness}) terbalik dibanding hasil perhitungan CRC32 di \textit{firmware} &
Membalik urutan \textit{byte} nilai \textit{checksum} menggunakan \texttt{.swap\_bytes()}
sebelum dibandingkan \\
\hline
2 &
\textit{Device offline} (\texttt{BadUserNameOrPassword}) setelah OTA berhasil
dan \textit{reboot} &
Paket \textit{firmware} OTA yang di-\textit{build} menyertakan token akses MQTT yang
berbeda (kedaluwarsa) dari token yang sedang aktif di ThingsBoard, karena
\texttt{wifi\_config.txt} sudah berubah sejak \textit{firmware} tersebut dikompilasi &
Membangun ulang (\textit{rebuild}) paket \textit{firmware} OTA menggunakan
\texttt{wifi\_config.txt} yang terbaru sebelum diunggah sebagai paket OTA baru \\
\hline
3 &
\textit{Firmware} yang sama diunduh ulang terus-menerus setiap \textit{device} melakukan
\textit{reconnect} MQTT (\textit{infinite loop}) &
\textit{Firmware} tidak pernah membandingkan versi \textit{firmware} yang ditetapkan
(\textit{assigned}) dengan versi yang sedang berjalan -- setiap menerima \textit{attribute}
\texttt{fw\_title}/\texttt{fw\_version}, unduhan langsung dimulai tanpa syarat &
Menambahkan pengecekan kesamaan judul dan versi \textit{firmware} sebelum memulai proses
unduh; unduhan di-\textit{skip} apabila versi sudah sama \\
\hline
\end{tabular}
\end{table}

\section{Verifikasi \textit{Checksum} CRC32}
Nilai CRC32 dihitung secara manual dan independen di sisi \textit{firmware} atas seluruh
data yang diunduh, kemudian dibandingkan dengan nilai target dari ThingsBoard.
Contoh luaran \textit{log} yang menunjukkan proses ini berhasil:

\begin{lstlisting}[caption={Cuplikan log verifikasi CRC32 yang berhasil}]
INFO - OTA: chunk 151 diterima (1024 bytes, total 619520/619520 bytes)
INFO - OTA: download selesai. CRC32 kita: 0x571B2CAE,
       target dari ThingsBoard: 0x571B2CAE (COCOK)
INFO - OTA sukses, reboot...
\end{lstlisting}

\section{Verifikasi Token Kredensial}
Untuk memastikan paket \textit{firmware} yang diunggah ke ThingsBoard benar-benar
menggunakan kredensial (token akses MQTT) yang aktif, dilakukan pemeriksaan \textit{string}
biner yang tertanam pada berkas \texttt{.bin} hasil kompilasi sebelum diunggah,
serta verifikasi kesamaan \textit{byte-per-byte} antara berkas lokal dan berkas yang
tersimpan di ThingsBoard (dengan mengunduh kembali paket yang telah diunggah dan
membandingkannya). Langkah ini terbukti efektif mendeteksi kasus ketidaksesuaian
token pada iterasi pengujian sebelumnya.

\section{Verifikasi Penghentian \textit{Loop} Unduhan}
Setelah perbaikan pada Tabel~\ref{tab:kendala} poin 3 diterapkan, pengujian ulang
menunjukkan bahwa \textit{device} berhenti mengunduh ulang \textit{firmware} yang versinya sama:

\begin{lstlisting}[caption={Log setelah perbaikan -- unduhan otomatis di-skip}]
INFO - MQTT connected ke ThingsBoard
INFO - OTA: coffee-enose v2.1.2 sudah versi yang jalan sekarang,
       skip download
\end{lstlisting}

\section{Hasil Akhir Pengujian}
Setelah ketiga perbaikan pada Tabel~\ref{tab:kendala} diterapkan, mekanisme OTA
terverifikasi berjalan secara utuh (\textit{end-to-end}) pada perangkat keras (\textit{hardware})
sebagai berikut: penetapan (\textit{assignment}) \textit{firmware} baru di ThingsBoard
$\rightarrow$ \textit{device} mendeteksi perbedaan versi $\rightarrow$ unduhan berjalan
per-\textit{chunk} (151 \textit{chunk}, total 619.520 \textit{byte}) $\rightarrow$ verifikasi
CRC32 cocok $\rightarrow$ \textit{device} melakukan \textit{restart} $\rightarrow$ \textit{device}
kembali terhubung ke MQTT dengan kredensial yang tetap valid $\rightarrow$ pada
percobaan koneksi berikutnya, unduhan tidak terulang karena versi telah sesuai.

Tampilan antarmuka \textit{dashboard} ThingsBoard \textit{Cloud} saat pemantauan telemetri dan status OTA \textit{device} dapat dilihat pada Gambar~\ref{fig:dashboard_thingsboard}.

\begin{figure}[H]
    \centering
    \includegraphics[width=0.8\textwidth]{thingboard.png}
    \caption{Tampilan \textit{Dashboard} ThingsBoard \textit{Cloud} untuk pemantauan data telemetri dan status OTA \textit{firmware update}.}
    \label{fig:dashboard_thingsboard}
\end{figure}

% =======================================================
\chapter{Kendala dan Pembelajaran}

\section{Konvensi \textit{Byte Order} yang Tidak Terdokumentasi}
Salah satu pembelajaran penting dari proyek ini adalah pentingnya melakukan
verifikasi empiris terhadap asumsi format data dari pihak ketiga (dalam hal ini,
representasi \textit{byte} dari nilai \textit{checksum} yang dikirim ThingsBoard), karena
perbedaan konvensi semacam ini tidak selalu terdokumentasikan secara eksplisit pada
dokumentasi resmi.

\section{Konsistensi Kredensial pada Proses \textit{Build}}
Karena kredensial WiFi/MQTT ditanamkan langsung ke dalam biner \textit{firmware} saat proses
kompilasi (\textit{compile-time embedding}), setiap kali paket \textit{firmware} baru akan
diunggah sebagai paket OTA, perlu dipastikan bahwa proses \textit{build} dilakukan
menggunakan berkas konfigurasi kredensial (\texttt{wifi\_config.txt}) yang paling
mutakhir, untuk menghindari \textit{device} kehilangan konektivitas setelah pembaruan.

\section{Pentingnya \textit{Guard} Perbandingan Versi}
Tanpa mekanisme pembanding versi yang eksplisit, sistem OTA rentan terhadap kondisi
\textit{infinite loop} yang tidak perlu, membebani \textit{bandwidth} dan siklus tulis \textit{flash}
memori \textit{device} secara berulang. Penambahan pengecekan versi sederhana terbukti cukup
untuk mencegah kondisi ini tanpa mengorbankan fungsi utama OTA.

% =======================================================
\chapter{Kesimpulan dan Saran}

\section{Kesimpulan}
\begin{enumerate}[label=\arabic*.]
    \item \textit{Firmware} akuisisi data sensor gas \textit{array} (8 kanal via 2x ADS1115) dan
    sensor lingkungan DHT22 pada ESP32-S3 berhasil diimplementasikan menggunakan
    Rust dengan \textit{framework} \texttt{esp-hal} dan \texttt{embassy}.
    \item Konektivitas WiFi dan MQTT ke ThingsBoard \textit{Cloud} berhasil diimplementasikan
    dengan mekanisme \textit{reconnect} otomatis dan pengiriman telemetri berkala.
    \item Mekanisme OTA \textit{firmware update} berhasil diimplementasikan dan diverifikasi
    berfungsi secara \textit{end-to-end} pada perangkat keras (\textit{hardware}) nyata, setelah melalui
    proses \textit{debugging} tiga kendala teknis signifikan: perbedaan urutan \textit{byte
    checksum} CRC32, inkonsistensi token kredensial pada paket OTA, dan
    ketiadaan pengecekan versi yang menyebabkan \textit{infinite loop} unduhan.
    \item Sistem OTA yang dihasilkan bersifat aman terhadap kegagalan
    (\textit{fail-safe}): apabila proses unduh gagal di tengah jalan, partisi
    \textit{firmware} yang sedang aktif tidak tersentuh, sehingga \textit{device} tetap dapat
    beroperasi dengan \textit{firmware} sebelumnya.
\end{enumerate}

\section{Saran}
\begin{enumerate}[label=\arabic*.]
    \item Melakukan pemasangan fisik seluruh sensor (2x ADS1115 dan DHT22) ke
    perangkat keras dan memvalidasi ulang seluruh jalur akuisisi data sebelum
    pengambilan data sampel kopi yang sesungguhnya dimulai.
    \item Melakukan kalibrasi masing-masing kanal sensor gas terhadap sampel
    referensi yang diketahui sebelum data digunakan untuk pelatihan model
    klasifikasi.
    \item Menambahkan mekanisme \textit{rollback} otomatis pada \textit{firmware} apabila
    \textit{firmware} hasil OTA gagal terhubung ke jaringan setelah sejumlah percobaan
    tertentu, sebagai lapisan keamanan tambahan di luar verifikasi \textit{checksum} CRC32.
    \item Mendokumentasikan prosedur pembaruan OTA (termasuk langkah verifikasi
    \textit{checksum} dan token) sebagai panduan baku bagi anggota tim lain yang
    akan melakukan pembaruan \textit{firmware} di kemudian hari.
\end{enumerate}

% =======================================================
\chapter*{Daftar Pustaka}
\addcontentsline{toc}{chapter}{Daftar Pustaka}

\begin{enumerate}[label={[\arabic*]}]
    \item Espressif Systems, \textit{esp-hal: Hardware Abstraction Layer for
    Espressif Chips}, \url{https://github.com/esp-rs/esp-hal}.
    \item Embassy Project, \textit{embassy: Async Embedded Framework for Rust},
    \url{https://embassy.dev}.
    \item ThingsBoard, \textit{ThingsBoard OTA Updates Documentation},
    \url{https://thingsboard.io/docs/user-guide/ota-updates/}.
    \item esp-hal-ota crate documentation, \url{https://crates.io/crates/esp-hal-ota}.
    \item OASIS, \textit{MQTT Version 3.1.1 Protocol Specification}.
    \item Texas Instruments, \textit{ADS1115 16-Bit ADC Datasheet}.
    \item Aosong Electronics, \textit{DHT22 (AM2302) Digital Temperature and
    Humidity Sensor Datasheet}.
    \item Repositori GitHub Proyek Coffee \textit{E-Nose} IoT, \url{https://github.com/yusuffebriii/Coffee-Enose-IoT-Project}.
\end{enumerate}

\end{document}
